\chapter{DP Optimizations}\label{ch:dpopt}

A correct DP that is too slow can usually be accelerated without changing its states, because the transitions that make it slow
(``the minimum over a range of earlier states'', ``the best split point'', ``the best earlier line'') have structure that a data
structure or a monotonicity argument can exploit. This chapter goes through the standard accelerations, roughly in increasing order of
difficulty, each with the shape of transition it recognises, the reason it is correct, and a problem solved with it.
Table~\ref{tab:dpopt} is the index; the rest of the chapter explains every row.

\begin{chapterintro}{\gold\ \plat\ \ioi}{Chapters~\ref{ch:dp} and~\ref{ch:dppatterns}, and the monotonic deque of
Section~\ref{sec:sliding}.}
  \item replace range transitions by prefix sums and monotone queues, including an exact $\Oh{nW}$ bounded knapsack;
  \item use the monotonicity of optimal split points: divide and conquer, monotone minima and SMAWK, and Knuth's optimization;
  \item turn ``the best earlier state'' into ``the lowest line'' with the convex hull trick, in its three forms;
  \item keep convex value functions by their breakpoints (slope trick), and remove a ``use exactly $k$'' dimension (Aliens trick);
  \item convolve convex sequences in linear time, and know when bitsets or swapping dimensions are the right tool.
\end{chapterintro}

\begin{table}[htbp]
\centering\small
\begin{tabularx}{\linewidth}{@{}>{\raggedright\arraybackslash}p{6.2cm}>{\raggedright\arraybackslash}p{3.4cm}X@{}}
\toprule
transition & technique & speed-up \\
\midrule
$dp_i = \sum_{j \in \text{window}} dp_j$, or the minimum over a sliding window & prefix sums, monotone queue & $\Oh{nk} \to \Oh{n}$ \\
$dp[g][i] = \min_j dp[g-1][j] + \mathit{cost}(j, i)$, best $j$ monotone in $i$ & divide and conquer & $\Oh{n^2} \to \Oh{n \log n}$ per layer \\
row minima of a totally monotone matrix & SMAWK & $\Oh{n \log n} \to \Oh{n}$ \\
$dp[l][r] = \min_m dp[l][m] + dp[m+1][r] + \mathit{cost}(l, r)$ & Knuth & $\Oh{n^3} \to \Oh{n^2}$ \\
$dp_i = \min_j (k_j x_i + m_j)$ & convex hull trick & $\Oh{n^2} \to \Oh{n}$ or $\Oh{n \log n}$ \\
a convex piecewise-linear value function & slope trick & $\Oh{nV} \to \Oh{n \log n}$ \\
``exactly $k$'' with an answer concave in $k$ & Aliens trick & removes the factor $k$ \\
$c_k = \min_{i + j = k} a_i + b_j$, $a$ and $b$ convex & merging slopes & $\Oh{nm} \to \Oh{n + m}$ \\
boolean rows shifted and OR-ed & bitset & divides by 64 \\
\bottomrule
\end{tabularx}
\caption{Which optimization fits which shape of transition.}
\label{tab:dpopt}
\end{table}

% ======================================================================================================================
\section[Windows and monotone queues]{Windows and monotone queues \hfill \gold}\label{sec:dpopt-window}

The simplest accelerations replace a loop over a range of earlier states. If $dp_i = \sum_{j=i-k}^{i-1} dp_j$, keep the prefix sums
of $dp$ and each state costs $\Oh{1}$. If $dp_i = c_i + \min_{i - k \le j < i} dp_j$, keep the monotonic deque of
Section~\ref{sec:sliding}: indices with increasing values, the front being the minimum of the window. Both turn $\Oh{nk}$ into
$\Oh{n}$.

\paragraph{The exact bounded knapsack.} Section~\ref{sec:knapsack} handled an item with $k$ copies by binary splitting, in
$\Oh{W \log k}$. A monotone queue does it in $\Oh{W}$. Fix the remainder $r$ of the capacity modulo the item's weight $w$: the
capacities $c = r, r + w, r + 2w, \dots$ form a chain, and with $t$ copies the new value at $c = r + tw$ is
\[
  \mathit{new}[r + tw] = \max_{t - k \le s \le t} \bigl(\mathit{dp}[r + sw] + (t - s)v\bigr)
  = tv + \max_{t - k \le s \le t} \bigl(\mathit{dp}[r + sw] - sv\bigr),
\]
the maximum of a sliding window of length $k + 1$ along the chain (Figure~\ref{fig:bknapmq}). Every capacity enters and leaves the
queue once per item, so each item costs $\Oh{W}$.

\begin{figure}[htbp]
\centering
\begin{tikzpicture}
  \node[blbl, anchor=east] at (-0.2, 0) {position $t$ on the chain};
  \arrs{C}{0}{0}{0/n, {}/dim, {}/dim, 1/cool, {}/dim, {}/dim, 2/cool, {}/dim, {}/dim, 3/hot, {}/dim, {}/dim, 4/n}
  \foreach \ua/\ub in {1/4, 4/7, 7/10, 10/13} \draw[->, line width=0.8pt, draw=fgline] (C-\ua.north) to[bend left=35] (C-\ub.north);
  \draw[mbrace] (C-4.south west) -- node[below=6pt, lbl] {window of $k + 1 = 3$ positions of the chain, $s = t - 2, \dots, t$}
    (C-10.south east);
  \node[notel, anchor=north west, text width=12cm] at (0, -1.25) {The cells are capacities $0, 1, 2, \dots$; for an item of
    weight $w = 3$ only capacities with the same remainder modulo $w$ interact, and the numbered cells are the chain $c = r +
    tw$. Along it, the new value at $t = 3$ is the best of the three window cells $s = 1, 2, 3$, each corrected by $-sv$; the
    deque gives that maximum in amortised $\Oh{1}$.};
\end{tikzpicture}
\caption{The bounded knapsack along one remainder class, with $k = 2$ copies.}
\label{fig:bknapmq}
\end{figure}

\begin{code}
void solve() {
    int n, cap;
    cin >> n >> cap;
    vi dp(cap + 1, 0), nd(cap + 1), qs(cap + 1), qv(cap + 1);
    rep(i, n) {
        ll w, v, k;
        cin >> w >> v >> k;
        // capacities with the same remainder mod w form a chain c = r + t*w:
        // nd[c] = t*v + max over t - k <= s <= t of (dp[r + s*w] - s*v),
        // a sliding-window maximum (a deque kept in the arrays qs, qv)
        rep(r, min<ll>(w, cap + 1)) {
            int head = 0, tail = 0;  // the window, values decreasing
            for (ll t = 0, c = r; c <= cap; t++, c += w) {
                ll val = dp[c] - t * v;
                while (tail > head && qv[tail - 1] <= val) tail--;
                qs[tail] = t, qv[tail++] = val;
                if (qs[head] < t - k) head++;  // at most k copies of item i
                nd[c] = qv[head] + t * v;
            }
        }
        swap(dp, nd);
    }
    print(dp[cap]);
}
\end{code}

The deque lives in two arrays reused for every chain, because creating a \cc{std::deque} per chain would cost more than the DP itself
when the weights are large. With 100 kinds of 1000 copies each and $W = 10^5$, binary splitting takes 0.09\,s and the queue 0.06\,s;
the gap grows with $k$, and the queue's bound does not depend on $k$ at all.

% ======================================================================================================================
\section[Divide and conquer optimization]{Divide and conquer optimization \hfill \plat}\label{sec:dnc}

A \emph{layered} DP splits the first $i$ elements into $g$ groups:
\[ dp_g(i) = \min_{j < i} \bigl(dp_{g-1}(j) + \mathit{cost}(j, i)\bigr), \]
where $\mathit{cost}(j, i)$ is the cost of the group $a_j, \dots, a_{i-1}$. Computed directly, each layer costs $\Oh{n^2}$. Let
$\mathit{opt}(i)$ be the best $j$ (the leftmost one, if several tie). In many problems $\mathit{opt}(i) \le \mathit{opt}(i + 1)$: the last
group never starts earlier when one more element arrives. Then compute the middle row $i$ first by scanning all candidates; its
optimum bounds both halves, rows left of the middle only scan $j \le \mathit{opt}(i)$ and rows right of it only $j \ge
\mathit{opt}(i)$ (Figure~\ref{fig:dnc}). Every level of the recursion scans about $n$ candidates in total, so a layer costs $\Oh{n \log n}$.

\begin{figure}[htbp]
\centering
\begin{tikzpicture}
  \tikzset{mc/.style={draw=fgline!80, line width=0.4pt, minimum size=5.6mm, inner sep=0pt, anchor=center},
    dq0/.style={fill=fgverm!45}, dq1/.style={fill=fgorange!40}, dq2/.style={fill=fgsky!45}, dq3/.style={fill=fggreen!30}}
  \foreach \ii in {1, ..., 8} {
    \node[lbl, anchor=east] at (-0.42, -0.56*\ii) {$i = \ii$};
    \foreach \jj in {0, ..., 7} \node[mc, fill=fgpaper] (X\ii\jj) at (0.56*\jj, -0.56*\ii) {};
  }
  \foreach \jj in {0, ..., 7} \node[lbl] at (0.56*\jj, -0.05) {\jj};
  \node[lbl, anchor=east] at (-0.42, -0.05) {$j$:};
  % scanned cells, coloured by recursion depth; the argmin of each row is marked
  \foreach \ii/\jlist/\dep in {4/{0,1,2,3}/0, 2/{0,1}/1, 7/{2,3,4,5,6}/1, 1/{0}/2, 3/{1,2}/2, 6/{2,3,4,5}/2, 8/{5,6,7}/2,
      5/{2,3,4}/3}
    \foreach \jj in \jlist \node[mc, dq\dep] at (X\ii\jj) {};
  \foreach \ii/\jj in {1/0, 2/1, 3/2, 4/2, 5/3, 6/4, 7/5, 8/5} \fill[fgink] (X\ii\jj) circle (1.5pt);
  \legendat{(5.0, -0.6)}{dq0/row 4 first: all candidates, dq1/next: rows 2 and 7, dq2/then rows 1 and 3; 6 and 8,
    dq3/last: row 5}
  \node[notel, anchor=north west, text width=5.8cm] at (4.75, -2.5) {Two groups over $3, 1, 4, 1, 5, 9, 2, 6$: row $i$ is the
    first $i$ elements, column $j$ the start of the last group, and the dot the optimum. Row 4 finds $\mathit{opt} = 2$, so rows 1
    to 3 scan only $j \le 2$ and rows 5 to 8 only $j \ge 2$. The recursion evaluates 24 of the 36 candidates; for large $n$ the
    saving is a factor $n / \log n$.};
\end{tikzpicture}
\caption{Divide and conquer optimization: each row's optimum bounds the rows below and above it in the recursion.}
\label{fig:dnc}
\end{figure}

\begin{problem}{Subarray Squares (CSES)}
Split an array of $n \le 3000$ positive numbers into $k$ contiguous parts so that the sum of the squares of the parts' sums is as
small as possible.
\solution $k$ layers of the DP above with $\mathit{cost}(j, i) = (\mathit{pre}_i - \mathit{pre}_j)^2$, each by divide and conquer:
$\Oh{kn \log n}$, 0.12\,s for $n = k = 3000$.
\end{problem}

\begin{code}
void solve() {
    int n, k;
    cin >> n >> k;
    vi pre(n + 1, 0);
    rep(i, n) {
        ll x;
        cin >> x;
        pre[i + 1] = pre[i] + x;
    }
    auto cost = [&](int j, int i) {  // the group a[j, i)
        ll s = pre[i] - pre[j];
        return s * s;
    };
    // prev[i]: the best split of the first i elements into g - 1 groups;
    // cur[i]: into g groups
    vi prev(n + 1, INF), cur(n + 1);
    prev[0] = 0;
    // fill cur[lo, hi) knowing that their best last split lies in [olo, ohi]
    auto go = [&](auto&& self, int lo, int hi, int olo, int ohi) -> void {
        if (lo >= hi) return;
        int mid = (lo + hi) / 2, best = olo;
        cur[mid] = INF;
        rep(j, olo, min(mid - 1, ohi) + 1) {  // the last group is a[j, mid)
            if (prev[j] == INF) continue;
            ll c = prev[j] + cost(j, mid);
            if (c < cur[mid]) cur[mid] = c, best = j;
        }
        self(self, lo, mid, olo, best);
        self(self, mid + 1, hi, best, ohi);
    };
    rep(g, 1, k + 1) {
        go(go, 0, n + 1, 0, n);
        swap(prev, cur);
    }
    print(prev[n]);
}
\end{code}

\usetemplate{Divide \& conquer optimization}{Ctrl+Alt+Y, Q}{cpdnc}

\begin{fact}[When is the optimum monotone?]
A sufficient condition is the \emph{quadrangle inequality}: $\mathit{cost}(a, c) + \mathit{cost}(b, d) \le \mathit{cost}(a, d) +
\mathit{cost}(b, c)$ for all $a \le b \le c \le d$ (``crossing is cheaper than nesting''). It holds for $(\text{sum of the
group})^2$ with non-negative values, for the number of pairs inside a group, for the variance of a group and for the distance of the
group's elements to their median. Its consequence is that the matrix $M[i][j] = dp_{g-1}(j) + \mathit{cost}(j, i)$ is \emph{totally
monotone}, which is exactly what the next section exploits further.
\end{fact}

\begin{tip}
In a contest, checking monotonicity with a brute force on random small inputs (compute every $\mathit{opt}(i)$ in $\Oh{n^2}$ and assert
that it never decreases) is faster and safer than proving the quadrangle inequality.
\end{tip}

% ======================================================================================================================
\section[Monotone minima and SMAWK]{Monotone minima and SMAWK \hfill \ioi}\label{sec:smawk}

The divide and conquer of the previous section is an instance of a general problem. A matrix $M$ (usually implicit: its entries are
computed on demand) is \emph{monotone} if the column of the leftmost minimum of each row never decreases from one row to the next.
Finding all row minima of an $n \times m$ monotone matrix takes $\Oh{(n + m) \log n}$ evaluations with the same recursion:
middle row first, then the halves with restricted columns. The template exposes exactly that, with the matrix given as a function.

\usetemplate{Monotone minima}{Ctrl+Alt+U, M}{cpmonominima}

If the matrix is \emph{totally} monotone (every $2 \times 2$ submatrix is monotone, as for a Monge matrix), the \defn{SMAWK} algorithm
finds all row minima with only $\Oh{n + m}$ evaluations. It alternates two steps. \emph{Reduce}: when there are more columns than
rows, compare entries along a stack of candidate columns and discard columns that can contain no row minimum, until at most $n$
remain. \emph{Interpolate}: solve the problem for the odd-numbered rows recursively, then fill in each even row by scanning only the
columns between the minima of its two neighbouring rows. Both steps are linear, and the recursion halves the rows. SMAWK matters when
$\log n$ is the difference between passing and failing, for example in some layered DPs with millions of states, or in the
$(\min, +)$ convolution of Section~\ref{sec:minplus}; in most contests the $\Oh{(n + m) \log n}$ recursion is enough and much simpler to
get right.

% ======================================================================================================================
\section[Knuth's optimization]{Knuth's optimization \hfill \plat}\label{sec:knuth}

For interval DPs of the form
\[ dp[l][r] = \min_{l \le m < r} \bigl(dp[l][m] + dp[m+1][r]\bigr) + \mathit{cost}(l, r), \]
let $\mathit{opt}(l, r)$ be the best split point. When $\mathit{cost}$ satisfies the quadrangle inequality and is monotone under
inclusion, the optimum is sandwiched between its neighbours (Figure~\ref{fig:knuth}):
\[ \mathit{opt}(l, r - 1) \le \mathit{opt}(l, r) \le \mathit{opt}(l + 1, r). \]
Restricting the loop over $m$ to that range makes the whole DP $\Oh{n^2}$: for a fixed length, the ranges of consecutive $l$ overlap
only at their ends, so their total length telescopes to $\mathit{opt}(n - \mathit{len}, n - 1) - \mathit{opt}(0, \mathit{len} - 2) + n
= \Oh{n}$, and there are $n$ lengths.

\begin{figure}[htbp]
\centering
\begin{tikzpicture}
  \tikzset{kc/.style={draw=fgline, line width=0.5pt, minimum size=7.5mm, inner sep=0pt, font=\ttfamily\small, anchor=center}}
  \foreach \cc [count=\rr from 0] in {2, 3, 1, 4, 2, 3} \node[font=\sffamily\small\bfseries, text=fgblue] at (0.75*\rr, 0.7) {\cc};
  \node[lbl, anchor=east] at (-0.45, 0.7) {$a$};
  \foreach \rr/\row in {0/{0,0,0,1,2,2}, 1/{1,1,2,2,3}, 2/{2,2,3,3}, 3/{3,3,3}, 4/{4,4}, 5/{5}}
    \foreach \val [count=\cc from 0] in \row {
      \pgfmathtruncatemacro{\col}{\cc + \rr}
      \node[kc] (K\rr\col) at (0.75*\col, -0.75*\rr) {\val};
    }
  \foreach \rr in {0, ..., 5} \node[lbl, anchor=east] at (-0.45, -0.75*\rr) {$l = \rr$};
  \node[kc, hot] at (K05) {2};
  \node[kc, cool] at (K04) {2};
  \node[kc, cool] at (K15) {3};
  \node[notel, anchor=west, text width=6.0cm] at (4.9, -1.6) {$\mathit{opt}(l, r)$ for Knuth Division on $2, 3, 1, 4, 2, 3$
    (row $l$, column $r$). The whole array only tries the cuts $2 \le m \le 3$, the range given by its left neighbour
    $\mathit{opt}(0, 4) = 2$ and the cell below, $\mathit{opt}(1, 5) = 3$. The values never decrease along a row or down a
    column.};
\end{tikzpicture}
\caption{The table of optimal split points of an interval DP is monotone in both directions.}
\label{fig:knuth}
\end{figure}

\begin{problem}{Knuth Division (CSES)}
An array of $n \le 5000$ numbers is cut into single elements; each step cuts one piece into two, at a cost equal to the sum of the
piece. Minimise the total cost.
\solution The first cut of $a_l, \dots, a_r$ splits it after some $m$, and both parts are then solved independently:
$\mathit{cost}(l, r)$ is the sum of the segment, which satisfies Knuth's conditions. The table of $2.5 \cdot 10^7$ cells is stored as a
triangle, rows of decreasing length, and the optimum table in 32 bits, which keeps the memory near 150\,MB; the program runs in
0.29\,s.
\end{problem}

\begin{code}
void solve() {
    int n;
    cin >> n;
    vi pre(n + 1, 0);
    rep(i, n) {
        ll x;
        cin >> x;
        pre[i + 1] = pre[i] + x;
    }
    // row l stores the segments [l, r] at index r - l (a triangle):
    // D[l][r - l] is the cheapest way to cut a[l..r] into single elements,
    // O[l][r - l] the best first cut (it splits off [l, O] and [O + 1, r])
    vector<vi> D(n);
    vector<vector<int32_t>> O(n);
    rep(l, n) D[l].assign(n - l, 0), O[l].assign(n - l, l);
    rep(len, 2, n + 1) rep(l, n - len + 1) {
        int r = l + len - 1;
        // Knuth: opt(l, r - 1) <= opt(l, r) <= opt(l + 1, r)
        int lo = O[l][r - 1 - l], hi = min<int>(O[l + 1][r - l - 1], r - 1);
        ll best = INF;
        rep(m, lo, hi + 1) {
            ll c = D[l][m - l] + D[m + 1][r - m - 1];
            if (c < best) best = c, O[l][r - l] = m;
        }
        D[l][r - l] = best + pre[r + 1] - pre[l];
    }
    print(D[0][n - 1]);
}
\end{code}

\usetemplate{Knuth optimization}{Ctrl+Alt+Y, Shift+Q}{cpknuth}

\begin{recognize}
Optimal merging or splitting of a sequence where the cost of an operation is the total of the piece (merging piles of stones, cutting
a stick or an array, Huffman-like costs with a fixed order), and optimal binary search trees. $n$ up to a few thousand with an
$\Oh{n^3}$ interval DP is the signal.
\end{recognize}

% ======================================================================================================================
\section[Convex hull trick]{The convex hull trick \hfill \plat}\label{sec:cht}

When the transition has the form $dp_i = \min_{j < i} (k_j x_i + m_j)$ (usually after expanding a product such as $s_i f_j$ or a square
$(\mathit{pre}_i - \mathit{pre}_j)^2$), every earlier state $j$ is a \emph{line} $y = k_j x + m_j$, and state $i$ asks which line is
lowest at $x = x_i$. The minimum of a set of lines is a concave piecewise-linear function, the \emph{lower envelope}; lines that never
appear on it can be discarded forever (Figure~\ref{fig:cht}).

\begin{figure}[htbp]
\centering
\begin{tikzpicture}
\begin{axis}[cpplot, width=12.5cm, height=6.0cm, xmin=0, xmax=100, ymin=0, ymax=9000, xlabel={$t$ (strength of the monster)},
  ylabel={time}, ytick={0, 3000, 6000, 9000}, scaled y ticks=false, legend pos=north west]
  \addplot[fgline, line width=0.8pt, domain=0:100, samples=2] {90*x + 2000};
  \addplot[fgline, line width=0.8pt, domain=0:100, samples=2] {60*x + 3000};
  \addplot[fgline, line width=0.8pt, domain=0:100, samples=2] {20*x + 4000};
  \addplot[fgline, line width=0.8pt, domain=0:100, samples=2] {10*x + 4800};
  \addplot[fgline, line width=0.8pt, domain=0:100, samples=2] {100*x};
  \addplot[fgline, line width=0.8pt, domain=0:100, samples=2] {20*x + 3000};
  \addplot[fgorange!90!black, line width=2.2pt, domain=0:37.5, samples=2] {100*x};
  \addplot[fgorange!90!black, line width=2.2pt, domain=37.5:100, samples=2] {20*x + 3000};
  \addplot[only marks, mark=*, mark size=2.2pt, fgblue] coordinates {(20, 2000) (30, 3000) (50, 4000) (90, 4800)};
  \node[font=\sffamily\scriptsize, text=fgorange!90!black, anchor=north west] at (axis cs:55, 3500) {lower envelope};
  \node[font=\sffamily\scriptsize, text=fgblue, anchor=north] at (axis cs:90, 4550) {$dp = 4800$};
\end{axis}
\end{tikzpicture}
\caption{Monster Game~I on the sample with $x = 100$: one line per earlier kill, $y = f_j t + dp_j$. Each query takes the lowest line
at $t = s_i$ (blue points, using only the lines that existed at that time). Of the six lines, only $100t$ and $20t + 3000$ end up on the
final lower envelope (orange).}
\label{fig:cht}
\end{figure}

\paragraph{Monotone slopes: a deque.} If the lines arrive with decreasing slopes, each new line enters the envelope at the right end, and
the lines it makes useless are at that end too: a line $b$ between $a$ and $c$ is useless when $c$ overtakes $b$ no later than $b$
overtakes $a$. So the envelope is a stack, each line pushed and popped once. If the queries also come in increasing order, a pointer
moving right answers them in amortised $\Oh{1}$; otherwise binary search on the envelope answers each in $\Oh{\log n}$. The template
compares the intersection points as exact fractions, so it never overflows or loses precision.

\begin{problem}{Monster Game I (CSES)}
A game has $n \le 2 \cdot 10^5$ levels. On each level you may kill or escape the monster, except on the last level, where it must be
killed. Killing a monster of strength $s$ takes $s \cdot f$ time, where $f$ is your current skill factor (initially $x$), and killing
monster $i$ sets $f$ to $f_i$. The strengths never decrease and the factors $f_i$ never increase. Minimise the total time.
\solution $dp_i = \min_j (dp_j + s_i f_j)$ over the previous kill $j$ (or the start, with $f = x$ and time 0): line $j$ has slope $f_j$
and intercept $dp_j$, and is queried at $s_i$. Slopes arrive non-increasing and queries non-decreasing, so the monotone hull with its
pointer runs in $\Oh{n}$; 0.08\,s.
\end{problem}

\begin{code}
void solve() {  // MonoCHT (cpmonocht) is pasted above solve()
    int n;
    ll x;
    cin >> n >> x;
    vi s(n), f(n);
    for (ll& v : s) cin >> v;
    for (ll& v : f) cin >> v;
    // dp[i]: the least time to kill monster i = min over the previous kill j
    // of dp[j] + s[i] * f[j]; the start is a kill with skill x at time 0.
    // Each kill j is the line y = f[j] * t + dp[j], queried at t = s[i].
    MonoCHT h;
    h.add(x, 0);
    ll dp = 0;
    rep(i, n) {
        dp = h.query_inc(s[i]);  // s is non-decreasing
        h.add(f[i], dp);  // the slopes f are non-increasing
    }
    print(dp);  // the last monster must be killed
}
\end{code}

\usetemplate{Monotone convex hull trick}{Ctrl+Alt+Y, Z}{cpmonocht}

\paragraph{Any order: LineContainer and Li Chao trees.} Without the monotonicity, the envelope must accept lines anywhere. A
\emph{LineContainer} keeps the lines of the envelope in a balanced tree ordered by slope, each storing the $x$ where the next line takes
over; inserting a line removes its useless neighbours, and a query is a binary search over those breakpoints, both $\Oh{\log n}$. A
\emph{Li Chao tree} is a segment tree over the $x$-domain in which every node keeps the line that is best at its midpoint; inserting a
line walks down one root path, swapping lines where the newcomer wins at the midpoint, and a query takes the minimum along the path to
$x$. It is $\Oh{\log C}$ per operation for a domain of size $C$, simpler to get right than the LineContainer, and it also accepts
line \emph{segments} (in $\Oh{\log^2 C}$).

\begin{problem}{Monster Game II (CSES)}
The same game without the monotonicity of the strengths and the factors.
\solution The same lines, inserted into a LineContainer; it keeps an upper envelope, so every line is stored negated: 0.09\,s.
\end{problem}

\begin{code}
void solve() {  // LineContainer (cplinecontainer) is pasted above solve()
    int n;
    ll x;
    cin >> n >> x;
    vi s(n), f(n);
    for (ll& v : s) cin >> v;
    for (ll& v : f) cin >> v;
    // the lines of Monster Game I, now in any order; the container keeps an
    // upper envelope, so each line is stored negated to get a minimum
    LineContainer h;
    h.add(-x, 0);
    ll dp = 0;
    rep(i, n) {
        dp = -h.query(s[i]);
        h.add(-f[i], -dp);
    }
    print(dp);
}
\end{code}

\usetemplate{LineContainer / Li Chao tree}{Ctrl+Alt+Y, L / Shift+L}{cplinecontainer, cplichao}

\begin{recognize}
After expanding the cost, the transition reads $dp_i = \min_j (A_j \cdot B_i + C_j) + D_i$: one product of a term of $j$ and a term of
$i$, plus terms of only one of them. Squares of prefix sums, ``rate times amount'' and ``distance times weight'' are the usual sources.
Which variant to use: monotone slopes and queries, the deque with a pointer; monotone slopes only, the deque with binary search;
neither, the LineContainer, or a Li Chao tree when the queries are integers in a known range or the lines are segments.
\end{recognize}

% ======================================================================================================================
\section[Slope trick]{Slope trick \hfill \ioi}\label{sec:slopetrick}

Some DPs keep a whole function of a value: $f_i(v)$, the least cost of the first $i$ elements if the $i$-th one ends at value $v$. When
every operation keeps $f_i$ \emph{convex and piecewise linear with integer slopes}, store it by its breakpoints: the multiset of points
where the slope increases by 1, plus the minimum value. ``Add $|v - a|$'' inserts the breakpoint $a$ twice (the slope rises by 2
there), ``take the prefix minimum'' ($\min_{u \le v} f(u)$) deletes every breakpoint to the right of the minimum, and shifts are
lazy offsets. Each becomes a few heap operations (Figure~\ref{fig:slopetrick}).

\begin{figure}[htbp]
\centering
\begin{tikzpicture}
  \foreach \px/\vals/\tx [count=\ci] in {0/{3,2,1,0,1,2}/{$f_1(v) = |v - 3|$},
      5.4/{4,2,2,2,3,4}/{$f_2$: prefix min of $f_1$, plus $|v - 1|$},
      10.8/{6,3,2,3,4,5}/{$f_3$: prefix min of $f_2$, plus $|v - 2|$}} {
    \begin{scope}[xshift=\px cm]
      \draw[->, line width=0.5pt, fgink!70] (0, 0) -- (4.2, 0) node[right, lbl] {$v$};
      \draw[->, line width=0.5pt, fgink!70] (0, 0) -- (0, 2.4);
      \foreach \vv in {0, ..., 5} \draw[line width=0.4pt, fgink!60] (0.7*\vv, 0.06) -- (0.7*\vv, -0.06) node[below, lbl] {\vv};
      \foreach \yy [count=\kk from 0] in \vals \coordinate (P\ci\kk) at (0.7*\kk, 0.35*\yy);
      \draw[line width=1.5pt, draw=fgblue] (P\ci0) -- (P\ci1) -- (P\ci2) -- (P\ci3) -- (P\ci4) -- (P\ci5);
      \foreach \kk in {0, ..., 5} \fill[fgblue] (P\ci\kk) circle (1.4pt);
      \node[lbl, anchor=south] at (2.1, 2.4) {\tx};
    \end{scope}
  }
  \fill[fgverm] (P13) circle (2.4pt);
  \fill[fgverm] (P21) circle (2.4pt); \fill[fgverm] (P23) circle (2.4pt);
  \fill[fgverm] (P32) circle (2.4pt);
\end{tikzpicture}
\caption{Slope trick on $a = 3, 1, 2$ (make it non-decreasing at the least total change). Each panel is the cost as a function of
the value $v$ of the last element; red dots mark the breakpoints that matter. The answer is the minimum of $f_3$, which is 2 (for
example $1, 1, 2$).}
\label{fig:slopetrick}
\end{figure}

\begin{problem}{Increasing Array II (CSES)}
Change elements of an array ($n \le 2 \cdot 10^5$) by $\pm 1$ per move until it is non-decreasing. Minimise the number of moves.
\solution Only the decreasing part of $f$ matters for the prefix minimum, so keep a max-heap of its breakpoints. For each $a$, push
$a$ (adding $|v - a|$); if the largest breakpoint is now above $a$, the cheapest final value for this prefix lies above $a$, so pay
the difference, pop it and push $a$ once more. $\Oh{n \log n}$, 0.07\,s. For a \emph{strictly} increasing array, apply the same to
$a_i - i$.
\end{problem}

\begin{code}
void solve() {
    int n;
    cin >> n;
    // f(v): the least cost so far if the last element ends at value v. f is
    // convex and piecewise linear; the max-heap holds the points where its
    // slope rises by 1, so its top is where f stops decreasing
    priority_queue<ll> pq;
    ll cost = 0;
    rep(i, n) {
        ll a;
        cin >> a;
        pq.push(a);  // add |v - a|
        if (pq.top() > a) {  // the cheapest end value lies above a: lower it
            cost += pq.top() - a;
            pq.pop();
            pq.push(a);
        }
    }
    print(cost);
}
\end{code}

\usetemplate{Slope trick}{Ctrl+Alt+Y, Shift+S}{cpslope}

\begin{recognize}
A DP over a sequence whose state is a \emph{value} with a huge range, and whose cost is a sum of absolute values or other convex
functions: make an array monotone at minimum total change, move tokens so that no two share a position, keep consecutive values within
a distance. If a brute-force DP over small values shows a convex, piecewise-linear row with integer slopes, the slope trick applies.
\end{recognize}

% ======================================================================================================================
\section[Aliens trick]{Aliens trick: Lagrangian relaxation \hfill \ioi}\label{sec:aliens}

``Choose \emph{exactly} $k$ things to maximise the total'' often has a natural DP with an extra dimension for the number chosen, which
multiplies the cost by $k$. If the best total $g(k)$ is \emph{concave} in $k$ (each extra item gains no more than the previous one),
the dimension can be removed. Charge a penalty $\lambda$ for every chosen thing and solve the unconstrained problem ``maximise $g(c) -
\lambda c$ over all counts $c$'', which is the original DP without the count. The optimal $c$ is where a line of slope $\lambda$
touches the graph of $g$ from above (Figure~\ref{fig:aliens}), so it decreases as $\lambda$ grows; binary search $\lambda$ until the
optimal count is $k$, and the answer is $(\text{the penalised optimum}) + \lambda k$. The cost is the DP's cost times the number of
search steps, about $\log$ of the slope range.

\begin{figure}[htbp]
\centering
\begin{tikzpicture}
\begin{axis}[cpplot, width=10.5cm, height=5.6cm, xmin=-0.3, xmax=9.3, ymin=0, ymax=22, xlabel={$k$, the number of subarrays},
  ylabel={$g(k)$}, xtick={0, ..., 9}]
  \addplot[fgblue, line width=1.3pt, mark=*, mark size=1.8pt] coordinates {(0,0) (1,7) (2,12) (3,15) (4,17) (5,18) (6,17)
    (7,15) (8,12) (9,6)};
  \addplot[fgorange!90!black, line width=1.1pt, domain=-0.3:9.3, samples=2] {2*x + 9};
  \addplot[only marks, mark=*, mark size=2.6pt, fgverm] coordinates {(3,15) (4,17)};
  \node[font=\sffamily\scriptsize, text=fgorange!90!black, anchor=east] at (axis cs:4.4, 20.6) {slope $\lambda = 2$};
\end{axis}
\end{tikzpicture}
\caption{The best total $g(k)$ of exactly $k$ disjoint subarrays of $5, -2, 4, -6, 3, -1, 2, -3, 4$ is concave. With a penalty of 2 per
subarray the penalised optimum is $g(k) - 2k = 9$, attained at $k = 3$ \emph{and} $k = 4$ (red): a tie, because the slope from 3 to 4
is exactly 2.}
\label{fig:aliens}
\end{figure}

\begin{problem}{Exactly $k$ subarrays}
Given $n \le 2 \cdot 10^5$ integers in $[-10^9, 10^9]$ and $k \le n$, choose exactly $k$ non-empty, pairwise disjoint subarrays with the
largest total sum.
\solution The total is concave in $k$ (it is a min-cost flow in disguise). For a penalty $\lambda$, a two-state DP (the current
element is or is not the end of a chosen subarray) finds the best penalised score in $\Oh{n}$, breaking ties towards \emph{fewer}
subarrays. Binary search the smallest $\lambda$ whose optimum uses at most $k$ subarrays; at that $\lambda$, $k$ is among the optimal
counts, and the answer is the penalised score plus $\lambda k$. About 50 steps of $\Oh{n}$: 0.09\,s.
\end{problem}

\begin{code}
void solve() {
    int n, k;
    cin >> n >> k;
    vi a(n);
    for (ll& x : a) cin >> x;
    // the best (score, -count) when any number of subarrays may be taken and
    // each one costs lam; among equal scores, fewer subarrays win
    auto run = [&](ll lam) {
        // out: the last element is in no subarray; in: it ends a subarray
        pll out{0, 0}, in{-INF, 0};
        for (ll x : a) {
            pll done = max(out, in);  // the best with everything so far closed
            pll ext{in.first + x, in.second};  // extend the open subarray
            pll start{done.first + x - lam, done.second - 1};  // or open one
            in = max(ext, start);
            out = done;
        }
        return max(out, in);
    };
    // the fewest-subarray optimum shrinks as lam grows: find the smallest lam
    // where it uses at most k subarrays, then add the k penalties back. Every
    // slope of the answer in k lies in [-1e9, 2e14], so this range is enough
    // and keeps the scores inside run far from overflow
    ll lo = -3e9, hi = 3e14;
    while (lo < hi) {
        ll mid = lo + (hi - lo) / 2;
        if (-run(mid).second <= k) hi = mid;
        else lo = mid + 1;
    }
    print(run(lo).first + lo * k);
}
\end{code}

\begin{pitfall}[Ties, and the range of the penalty]
Several counts can be optimal for the same $\lambda$, as in Figure~\ref{fig:aliens}, so the DP may never report exactly $k$. Break ties
consistently (here: fewest subarrays), search for the smallest $\lambda$ with count $\le k$, and add back $\lambda k$ with the
\emph{requested} $k$, not the count the DP returned. Second, bound $\lambda$ by the possible slopes of $g$, not by something huge:
with $\lambda = -3 \cdot 10^{14}$ the DP above would reward $2 \cdot 10^5$ subarrays with $6 \cdot 10^{19}$, past the range of 64-bit
integers. The first version of this solution made exactly that mistake; small random tests passed, and only a full-size test showed
the overflow.
\end{pitfall}

% ======================================================================================================================
\section[Convex (min, +) convolution]{Convex $(\min, +)$ convolution \hfill \ioi}\label{sec:minplus}

The $(\min, +)$ convolution $c_k = \min_{i + j = k} (a_i + b_j)$ appears whenever two independent parts share a budget: $a_i$ is the
best cost of spending $i$ units on the first part, $b_j$ on the second. In general it costs $\Oh{nm}$, but two special cases are fast.
\begin{itemize}
  \item \textbf{Both convex}: the result is convex, and its sequence of slopes is the two slope sequences merged in sorted order, like
        merging two sorted lists (the Minkowski sum of two convex polygonal lines, Figure~\ref{fig:minplus}): $\Oh{n + m}$. The
        reason is that, with convex costs, the cheapest next unit is always the cheaper of the two parts' next units.
  \item \textbf{One convex, one arbitrary}: the matrix $M[k][j] = b_j + a_{k - j}$ is totally monotone when $a$ is convex, so its row
        minima come from monotone minima in $\Oh{(n + m) \log}$, or from SMAWK in $\Oh{n + m}$.
\end{itemize}

\begin{figure}[htbp]
\centering
\begin{tikzpicture}[x=0.85cm, y=0.55cm]
  \draw[line width=0.4pt, fgline] (0, 0) -- (3.3, 0);
  \draw[line width=1.5pt, draw=fgblue, mark=*, mark size=1.4pt, mark options={fill=fgblue}]
    plot coordinates {(0,0) (1,-3) (2,-4) (3,-2)};
  \node[lbl, anchor=south west] at (0, 0.3) {$a$: slopes $-3, -1, 2$};
  \begin{scope}[xshift=4.0cm]
    \draw[line width=0.4pt, fgline] (0, 0) -- (3.3, 0);
    \draw[line width=1.5pt, draw=fgorange!90!black, mark=*, mark size=1.4pt, mark options={fill=fgorange!90!black}]
      plot coordinates {(0,0) (1,-2) (2,-2) (3,-1)};
    \node[lbl, anchor=south west] at (0, 0.3) {$b$: slopes $-2, 0, 1$};
  \end{scope}
  \begin{scope}[xshift=8.4cm]
    \draw[line width=0.4pt, fgline] (0, 0) -- (6.3, 0);
    \draw[line width=1.5pt, draw=fggreen!60!black, mark=*, mark size=1.4pt, mark options={fill=fggreen!60!black}]
      plot coordinates {(0,0) (1,-3) (2,-5) (3,-6) (4,-6) (5,-5) (6,-3)};
    \node[lbl, anchor=south west] at (0, 0.3) {$c$: slopes $-3, -2, -1, 0, 1, 2$};
  \end{scope}
\end{tikzpicture}
\caption{Merging the slopes of two convex sequences gives their $(\min, +)$ convolution. Here the result is $c = (0, -3, -5, -6, -6,
-5, -3)$, exactly what the brute force $\min_{i + j = k}(a_i + b_j)$ computes.}
\label{fig:minplus}
\end{figure}

\usetemplate{$(\min, +)$ convolution (convex)}{Ctrl+Alt+U, Shift+C}{cpminplus}

A typical use is a knapsack with \emph{small weights} but many items: group the items by weight $w$; within a group, taking the $c$
most valuable items is best, and the values of the first $c$ items form a concave sequence. Merging a group into the table is then,
for every remainder modulo $w$, a $(\max, +)$ convolution of the table (arbitrary) with that concave sequence: monotone minima per
remainder class instead of a loop over all counts.

% ======================================================================================================================
\section[Bitsets and swapped dimensions]{Bitsets and swapped dimensions \hfill \gold}\label{sec:dpopt-bitset}

\paragraph{Bitsets.} A boolean DP whose transition shifts a whole row, like subset sum ($\mathit{can}[s] \mathrel{|}= \mathit{can}[s
- a_i]$), runs 64 times faster on a \cc{bitset}: \cc{b |= b << a} (Section~\ref{sec:bitset} measures a $14\times$ speed-up on Money
Sums). When the items sum to $S$, there are at most $\Oh{\sqrt S}$ distinct values among them, since $1 + 2 + \dots + t \le S$
forces $t \le \sqrt{2S}$. Group equal values and split each group into bundles of $1, 2, 4, \dots$ copies (Section~\ref{sec:knapsack}):
the bitset subset sum then costs $\Oh{S \sqrt S \log / 64}$ in the worst case, and much less in practice, the standard solution to
``can these be split into two equal halves'' with $S$ up to $10^6$.

\usetemplate{Subset sums with bitset}{Ctrl+Alt+T, Shift+1}{cpbitsetsum}

\paragraph{Swapping the value and a dimension.} When a state would include a huge quantity (a weight up to $10^9$) but another
quantity is small (at most 100 items, values summing to $10^5$), make the small quantity the index and the huge one the value that is
minimised: $dp[\mathit{value}]$ is the least weight that reaches it. The LIS tail array (Section~\ref{sec:lis}) is the same idea:
instead of ``the longest subsequence ending at each value'', ``the smallest last value for each length''. Asking ``which quantity is
small enough to be an index?'' is often the whole optimization.

\begin{summary}[Chapter summary]
\begin{itemize}
  \item Range transitions become $\Oh{1}$ with prefix sums or a monotone queue; along each remainder class, a queue gives the exact
        $\Oh{nW}$ bounded knapsack.
  \item When the optimal split point is monotone (quadrangle inequality), divide and conquer computes a layer in $\Oh{n \log n}$,
        SMAWK in $\Oh{n}$, and Knuth's sandwich $\mathit{opt}(l, r-1) \le \mathit{opt}(l, r) \le \mathit{opt}(l+1, r)$ makes an
        interval DP $\Oh{n^2}$.
  \item Products of a $j$-term and an $i$-term make each earlier state a line: a deque for monotone slopes, a LineContainer or a Li
        Chao tree otherwise.
  \item Convex piecewise-linear value functions are stored by their breakpoints in heaps (slope trick).
  \item ``Exactly $k$'' with a concave answer becomes a penalty $\lambda$ found by binary search; break ties consistently and bound
        $\lambda$ by the slopes.
  \item Convex sequences convolve by merging slopes; bitsets divide boolean DPs by 64.
\end{itemize}
\end{summary}

\begin{practice}
\begin{pplist}
  \item Subarray Squares
  \item Houses and Schools
  \item Knuth Division
  \item Monster Game I
  \item Monster Game II
  \item Increasing Array II
  \item Money Sums
  \item Book Shop
  \item Elevator Rides
\end{pplist}
\end{practice}
